\documentclass[border=14pt]{standalone}
\usepackage{ctex}
\usepackage{tikz}
\begin{document}
\begin{tikzpicture}[x=1cm, y=1cm, font=\normalsize, line width=0.8pt]

% fig-ge-05 三相对称分量的三张脸（中文）
% 三子图并排：正序（A→B→C 逆时针 120°）、负序（反向）、零序（三支同向重合）
% 禁令②：零序必须一眼看出「三支同向」——画三支平行略错开并注明同相位

\def\R{1.5}
\def\cw{4.4}
\def\gap{1.2}

% ── 1 正序 ──
\begin{scope}[shift={(0,0)}]
  \node[anchor=south] at ({\cw/2},3.8) {正序};
  \coordinate (O) at ({\cw/2},1.6);
  \draw[->] (O) -- ++(90:\R)   node[anchor=south east, font=\small] {A};
  \draw[->] (O) -- ++(-30:\R)  node[anchor=north west, font=\small] {B};
  \draw[->] (O) -- ++(210:\R)  node[anchor=north east, font=\small] {C};
  \draw[->, line width=0.5pt, gray] (O) ++(90:0.5) arc (90:-30:0.5);
  \node[anchor=north, font=\small, align=center] at ({\cw/2},-0.7) {A→B→C 顺时针 120°\\大小相等、互差 120°};
\end{scope}

% ── 2 负序 ──
\begin{scope}[shift={(5.6,0)}]
  \node[anchor=south] at ({\cw/2},3.8) {负序};
  \coordinate (O) at ({\cw/2},1.6);
  \draw[->] (O) -- ++(90:\R)   node[anchor=south east, font=\small] {A};
  \draw[->] (O) -- ++(210:\R)  node[anchor=north east, font=\small] {B};
  \draw[->] (O) -- ++(-30:\R)  node[anchor=north west, font=\small] {C};
  \draw[->, line width=0.5pt, gray] (O) ++(90:0.5) arc (90:210:0.5);
  \node[anchor=north, font=\small, align=center] at ({\cw/2},-0.7) {A→C→B 逆时针 120°\\大小相等、互差 120°};
\end{scope}

% ── 3 零序（三支同向，略错开画清）──
\begin{scope}[shift={(11.2,0)}]
  \node[anchor=south] at ({\cw/2},3.8) {零序};
  % 三支同向（都向上），横向错开 0.22 以便看清是三支
  \draw[->, line width=1pt] ({\cw/2-0.30},0.6) -- ++(90:\R)
    node[anchor=south east, font=\small] {A};
  \draw[->, line width=1pt] ({\cw/2},0.6) -- ++(90:\R)
    node[anchor=south, font=\small] {B};
  \draw[->, line width=1pt] ({\cw/2+0.30},0.6) -- ++(90:\R)
    node[anchor=south west, font=\small] {C};
  \node[anchor=north, font=\small, align=center] at ({\cw/2},-0.7) {三支同向、同相位\\大小相等};
\end{scope}

\end{tikzpicture}
\end{document}